% !TEX root = ../main.tex
%-----------------------------------------------------------------------
%  Summary and resource request.
%  \input by ../main.tex -- not compiled on its own.
%-----------------------------------------------------------------------
%-----------------------------------------------------------------------
\section*{Summary and resource request}
\addcontentsline{toc}{section}{Summary and resource request}
%-----------------------------------------------------------------------

\begin{center}
\begin{tabularx}{\textwidth}{@{}lL@{}}
\toprule
\textbf{Compute} & \textbf{$\sim$75 node-hours} expected;
  \textbf{200 requested} as a ceiling\\
\textbf{Node type} & CPU only, MPI, no GPU. 27 ranks per run\\
\textbf{Peak concurrency} & 2 nodes (4 runs packed per 128-core node), or 7 if
  run one per node\\
\textbf{Storage} & $<$ 100 GB working; $\sim$23 GB retained\\
\textbf{Duration} & 2--3 weeks including setup and analysis\\
\textbf{Software} & MITgcm with a Tapenade-generated adjoint\\
\bottomrule
\end{tabularx}
\end{center}

The request is revised down from an earlier estimate of 150 node-hours with a
400-hour ceiling. Both the forward and the adjoint cost have since been measured
directly from completed runs rather than estimated, and a configuration change
identified in the meantime removes roughly 60\,\% of the adjoint cost.
Section~\ref{sec:compute} gives the measurements and states what remains
uncertain.

\paragraph{What is being proposed.}
Seven adjoint runs of an idealised single-basin ocean, identical to a completed
reference run in every respect except the vertical diffusivity $\kv$, which is
varied over a factor of 128. The runs establish how strongly adjoint sensitivity
patterns depend on the mixing parameterisation, and they provide the first
designed training data for a neural-network surrogate of those sensitivities.

\paragraph{What is being asked of reviewers.}
Section~\ref{sec:questions} lists six open questions. Two of them --- whether ten
years is enough re-equilibration, and whether the perturbation should carry
spatial structure --- change the experiment if answered differently, and are best
settled before the pilot member runs. Section~\ref{sec:sequence} places a formal
review point after the pilot, at roughly 15\,\% of the total cost.


\noindent
Sections~\ref{sec:motivation} and~\ref{sec:modelconfig} are common to both
parts. Part~I is self-contained from Section~\ref{sec:objective} onwards and
carries the appendices; a reviewer of the resource request needs Part~I and the
appendices only.

